\section{Lagrangian contact structures}
\label{Lagrangian contact structures.}

Let us recall Takeuchi's paper \cite{Tak} for Lagrangian contact structures.
A contact structure with a Lagrangian pair is called a~Lagrangian contact structures in~\cite{Tak}.
A typical example of Lagrangian contact structures
is the projective cotangent vector bundle
$M = P(T^*W)$ of an $(n+1)$-dimensional manifold~$W$
with an af\/f\/ine structure (a~torsion-free linear connection) or
a~projective structure
(a~projective equivalence class of torsion-free linear connections).
For the canonical contact structure~$D$ on~$M$,
we take as a Lagrangian pair horizontal and vertical vector bundles
(cf.\ Example~\ref{hor+ver}).
In~\cite{Tak}, it is given the description of Cartan connections on~$P(T^*W)$
associated to Lagrangian contact structures and the equivalence of the vanishing
of the curvature of Cartan connection and the projective f\/latness of $W$.

The f\/lat model, which is a homogeneous space qualif\/ied as a model
for a Cartan connection, with Lagrangian contact structure is
the projective cotangent bundle $P(T^*P^{n+1})$
of the ($n+1$)-dimensional projective space $P^{n+1}=P(\mathbf{R}^{n+2})$.

Put $G={\rm PGL}(n+2,{\mathbf{R}})={\rm GL}(n+2,{\mathbf{R}})/C$ ($C$ is the center,
$C=\mathbf{R}^{\times}\cdot I_{n+2}$),
and $\mathfrak{g}=\operatorname{Lie} G \cong \mathfrak{sl}(n+2,{\mathbf{R}})$.
The Lie algebra $\mathfrak{g}$ has a structure of
a simple graded Lie algebra (GLA) of second kind as follows:
\begin{gather*}
       {\mathfrak g} = {\mathfrak{sl}}(n+2,{\mathbf{R}})  = {\mathfrak g}_{-2}\oplus {\mathfrak g}_{-1} \oplus
                    {\mathfrak g}_0
                    \oplus {\mathfrak g}_1 \oplus {\mathfrak g}_2 \\
\hphantom{{\mathfrak g}}{}
= \left\{{\begin{pmatrix}
                       0 & 0 & 0 \\
                       0 & \mathrm{O}_n & 0 \\
                       a & 0 & 0
                    \end{pmatrix}}\right\}
                       \oplus
                    \left\{{\begin{pmatrix}
                       0 & 0 & 0 \\
                       b_1 & \mathrm{O}_n & 0 \\
                       0 & {}^tb_2 & 0
                    \end{pmatrix}}\right\}
                       \oplus
                    \left\{{\begin{pmatrix}
                       \alpha & 0 & 0 \\
                       0 & A & 0 \\
                       0 & 0 &\beta
                    \end{pmatrix}}\right\} \\
\hphantom{{\mathfrak g}=}{} \quad {}\oplus
                    \left\{{\begin{pmatrix}
                       0 & {}^tc_1 & 0 \\
                       0 & \mathrm{O}_n & c_2 \\
                       0 & 0 & 0
                    \end{pmatrix}}\right\}
                       \oplus
                    \left\{{\begin{pmatrix}
                       0 & 0 & d \\
                       0 & \mathrm{O}_n & 0 \\
                       0 & 0 & 0
                     \end{pmatrix}}\right\}, \\
  (a,d,\alpha,\beta \in {\mathbf R},\,  b_1,b_2,c_1,c_2\in
                    {\mathbf R}^n, \, A\in \mathfrak{gl}(n,{\mathbf{R}}); \,
                    \alpha +\beta +\operatorname{tr}A=0),
                  \qquad [\mathfrak{g}_p,\mathfrak{g}_q] \subset
                    \mathfrak{g}_{p+q}.
\end{gather*}
Put
\begin{gather*}
\mathfrak{m}= \mathfrak{g}_{-2}\oplus \mathfrak{g}_{-1},
\end{gather*}
then $\mathfrak{m}$ is a fundamental GLA of contact type, i.e., Heisenberg
algebra.
Put
\begin{gather*}
\mathfrak{g}'=\mathfrak{g}_0 \oplus \mathfrak{g}_1 \oplus
\mathfrak{g}_2.
\end{gather*}
Let $G'$ be the Lie subgroup of $G = {\rm PGL}(n+2, {\mathbf{R}})$ def\/ined by
\begin{gather*}
G' = P\left\{
\begin{pmatrix}
* & * & * \\
0 & * & * \\
0 & 0 & *
\end{pmatrix}
\in
{\rm GL}(n+2, {\mathbf{R}})
\right\}.
\end{gather*}
Then the Lie algebra of $G'$ is given by $\mathfrak{g}'$. Note that
$\dim \mathfrak{g}' = n^2 + 2n + 2$.

The group $G$ transitively acts on the f\/lag manifold
\begin{gather*}
P\big(T^*P^{n+1}\big) = \big\{ V_1 \subset V_{n+1} \subset {\mathbf{R}}^{n+2} \,|\,
\dim V_1 = 1, \dim V_{n+1} = n + 1 \big\} \subset P^{n+1}\times P^{n+1*}.
\end{gather*}
Then $G'$ is the isotropy group of $(V_1, V_{n+1}) = (\langle e_0 \rangle,
\langle e_0, \dots, e_n \rangle)$ for the standard basis $e_0, e_1, \dots$,
$e_n, e_{n+1}$ of ${\mathbf{R}}^{n+1}$.
Therefore we have
\begin{gather*}
G/G' \cong P\big(T^*P^{n+1}\big) \qquad \big(\cong P\big(T^*P^{n+1*}\big)\big).
\end{gather*}
Note that $\mathfrak{g}_{-1} \subset \mathfrak{m} = T_o(G/G')$, where
$o = G'$ is the origin,
def\/ines the contact structure $D$ on $G/G'$ which corresponds to
the canonical contact structure
on $P(T^*P^{n+1})$ via the above dif\/feomorphism (cf.~\cite{I-M}).

Next, we consider
\begin{gather*}  \mathfrak{e}^1=\left\{{\begin{pmatrix}
                      0 & 0 & 0 \\
                      b_1 & \mathrm{O}_n & 0 \\
                      0 & 0 & 0
                    \end{pmatrix}}\right\}, \;
  \mathfrak{e}^2=\left\{{\begin{pmatrix}
                      0 & 0 & 0 \\
                      0 & \mathrm{O}_n & 0 \\
                      0 & {}^tb_2 & 0
                     \end{pmatrix}}\right\}.
\end{gather*}
Then we have
\begin{gather*}
\mathfrak{g}_{-1}=\mathfrak{e}^1\oplus \mathfrak{e}^2,\qquad
[\mathfrak{e}^1,\mathfrak{e}^1]=[\mathfrak{e}^2,\mathfrak{e}^2]=0,\qquad
\mathfrak{g}_{-2}=[\mathfrak{e}^1,\mathfrak{e}^2],
\\
[\mathfrak{g}_0,\mathfrak{e}^1]
       \subset \mathfrak{e}^1,\qquad
       [\mathfrak{g}_0,\mathfrak{e}^2]
       \subset \mathfrak{e}^2.
\end{gather*}
Denote by $E_{ij}\in \mathfrak{gl}(n+2,{\mathbf{R}})$ the matrix unit of
$(i,j)$ component. Then we put
\begin{gather*}
\gamma =E_{n+2,1}\in \mathfrak{g}_{-2},\\
e_i =E_{i+1,1}\in
\mathfrak{e}^1\subset \mathfrak{g}_{-1},\qquad f_i=E_{n+2,i+1}\in
\mathfrak{e}^2\subset \mathfrak{g}_{-1}\qquad (1\leq i\leq n).
\end{gather*}
We set
\begin{gather*}
[X,Y] = - A(X,Y)\gamma \qquad (X,Y\in \mathfrak{g}_{-1}).
\end{gather*}
It follows that
\begin{gather*}
A(e_i,f_j)={\delta}_{ij}, \qquad A(e_i, e_j) = 0, \qquad A(f_i, f_j) = 0 \qquad  (1 \leq i, j \leq n).
\end{gather*}
Therefore we have that $A$ is a symplectic form.
Then $\mathfrak{g}_{-1}$ becomes a symplectic vector space with respect to~$A$,  and
$e_1, \dots, e_n$, $f_1, \dots, f_n$ form a symplectic basis of the symplectic vector space~$(\mathfrak{g}_{-1},A)$.
Moreover
$\mathfrak{e}^1$, $\mathfrak{e}^2$ form a Lagrangian pair of~$(\mathfrak{g}_{-1}, A)$.

Moreover, put
\begin{gather*}
\mathfrak{a}^1=\mathfrak{e}^1+\mathfrak{g}', \qquad
\mathfrak{a}^2=\mathfrak{e}^2+\mathfrak{g}'.
\end{gather*}
Then we easily verify that
\begin{gather*}
\mathrm{Ad}(G')\mathfrak{a}^1=\mathfrak{a}^1, \qquad
\mathrm{Ad}(G')\mathfrak{a}^2=\mathfrak{a}^2,\qquad
\mathfrak{a}^1\cap \mathfrak{a}^2=\mathfrak{g}'.
\end{gather*}
Thus $\mathfrak{a}^1$ and $\mathfrak{a}^2$ induce invariant
dif\/ferential systems~$E_1$ and $E_2$ on $G/G'$ respectively.
The pair $(E_1,E_2)$ forms
a Lagrangian pair of $(G/G', D)$ and thus
that of the standard contact structure $D$ on $P(T^*P^{n+1})$.
Moreover
both $E_1$ and $E_2$ are completely integrable.
It follows that
$P(T^*P^{n+1})/\mathcal{E}_2\cong P^{n+1}$, $P(T^*P^{n+1})/\mathcal{E}_1
\cong {P^{n+1}}^*$,
where
$\mathcal{E}_1$, $\mathcal{E}_2$ are the foliations induced by~$E_1$,~$E_2$
respectively.

For the linear isotropy representation
$\rho \colon  G' \longrightarrow \mathrm{GL}(\mathfrak{m})$ at the origin
$o = G'$ in $G/G'$,
we have, with respect to the basis $\{ \gamma, e_1, \dots, e_n, f_1, \dots, f_n \}$ in $\mathfrak{m}$,
\begin{gather*}
\tilde{G} =\rho(G')
          =\left\{{\begin{pmatrix}
                      a & 0 & 0 \\
                      b_1 & A & \mathrm{O}_n \\
                      b_2 & \mathrm{O}_n & a\, {}^t\! A^{-1}
                    \end{pmatrix}} \Bigg\vert \,
a\in \mathbf{R}^*,\, A\in \mathrm{GL}(n,{\mathbf{R}}), \,b_1, b_2\in \mathbf{R}^n
\right\}.
\end{gather*}

Then the $\tilde{G}$-structures of type $\mathfrak{m}$ are in bijective
correspondence with the Lagrangian contact structures~\cite[Theorem~5.1]{Tak}.
Note that $\tilde{G}$-structure is of inf\/inite type~\cite[Chapter~I]{K}.
How\-ever~${\mathfrak g}$ is the prolongation of $({\mathfrak m}, {\mathfrak g}_0)$
in the sense of Tanaka~\cite{Tan1,Y}. Moreover
$\tilde{G} = G_0^{\#}$ in the notation of~\cite{Tan}.
